\documentclass[11pt,a4paper]{article}
\usepackage[T1]{fontenc}
\usepackage[utf8]{inputenc}
\usepackage{amsmath,amssymb,xcolor}
\usepackage[margin=25mm]{geometry}
\usepackage[colorlinks=true,urlcolor=blue]{hyperref}
\definecolor{nexorablue}{HTML}{2459B5}
\setlength{\parindent}{0pt}
\setlength{\parskip}{7pt}
\setlength{\emergencystretch}{2em}
\renewcommand{\familydefault}{\sfdefault}
\begin{document}

{\small\bfseries\color{nexorablue} NEXORA CPGE}\par
{\small MP / NEX-STR-01 / Corrigé}\par
{\LARGE\bfseries Un sous-groupe défini par une contrainte additive}\par
\vspace{0.25em}\hrule\vspace{0.6em}

\section*{Sous-groupe}
Le couple nul appartient à $H$. Pour deux couples de $H$, la différence de leurs expressions $a-2b$ est encore divisible par $3$. La différence des couples appartient donc à $H$.

\section*{Générateurs}
La relation $a=3k+2b$ donne $(a,b)=k(3,0)+b(2,1)$. Réciproquement, une telle combinaison vérifie le critère. La seconde coordonnée détermine $b$, puis la première détermine $k$.


\end{document}
